\section{验证结果与完成度}
\label{sec:evaluation}
\subsection{评测方法与证据分层}
本地材料包含不同阶段、模型和执行模式的测试，不能合并为一个总体通过率。离线固定消息用于验证程序路径；真实模型预览用于检查语义和回答；原子 MCP 测试检查具体服务；模拟入站用于验证事件转换与参与；QQ 投递则需独立查看输出记录。

本报告引用已存资料时保留原测试日期和限制，并在报告编制过程中重新执行一轮不联网的 100 条固定消息流程。未为撰写报告重新向 QQ 群发送消息，也未调用付费模型生成新测试答案。验证中涉及的真实查询结果均来自既有记录。

\begin{longtable}{@{}L{26mm}L{61mm}L{62mm}@{}}
\caption{主要证据及其适用范围}\\\toprule
证据 & 输入与执行范围 & 能够支持的结论\\\midrule\endfirsthead
\toprule 证据 & 输入与执行范围 & 能够支持的结论\\\midrule\endhead
100 条固定消息 & 脚本模型、校园 fixture、内存输出，经过生产组合；真实 QQ 发送为 0 & 验证路径、参数投影、调用次数、结果状态及去重；不衡量真实模型回答准确率。\\
26 项 MCP 查询 & 录制环境真实服务的原子调用；包含可选服务案例 & 26 个选定案例通过，证明相应服务可调用；不代表全部自然语言任务。\\
5 项自然语言问答 & 真实 DeepSeek Flash、真实业务资料、Runtime 预览 & 代表问题调用预期能力并得到已核对的结果；样本较少且非随机抽样。\\
四人 20 条讨论 & 固定台词、模拟发言者、正式参与路径、真实模型与评课服务 & 代表场景能从关闭评课到授权启用，再查询并回应群成员需求。\\
原生入站排练 & OneBot 转换与正式接收路径；记录 1 次 output call 和 canary 完成 & 已存记录支持一次回复路径完成；模拟成员输入不能称为四个真实用户账号完整发言。\\
广覆盖群聊预览 & 93 个真实模型预览请求，另 7 项隔离测试 & 暴露语义、格式、来源和上下文问题；保留部分通过与失败结果。\\
Web 与组件回归 & 指定版本的自动化和浏览器场景 & 支持相应回归行为；不能证明所有线上服务配置、所有真实群或所有设备。\\\bottomrule
\end{longtable}

\subsection{离线 100 条固定消息流程}
既有 2026 年 9 月 6 日录制结果为 94 条 Runtime 完成响应，6 条由入口规则保留为兼容路径处理。执行统计为 187 次脚本模型调用、53 次 MCP fixture 调用和 11 次本地校车能力调用；真实 QQ 发送、Memory 写入均为 0。

动作、结果、投递预期、工具名、能力步骤和重复重放检查均无不匹配。每条业务计划不超过一步。这里的“兼容路径”是测试中记录的入口处置，不表示旧 1.0 智能体在生产中继续处理这些请求。

编制报告时按第\ref{sec:deployment}节命令重新运行同一 fixture，结果见随源码提供的离线基准摘要与检查记录。这轮复测用于校验报告中的流程统计与当前工作区一致，不增加真实模型质量样本。

\subsection{真实模型代表任务}
以下数据来自已存 \code{manual-answers.json}，运行于 2026 年 9 月 6 日录制验证。使用 Flash 标准档、关闭额外 thinking，实际读取业务资料。耗时是单次观测，包含整体请求过程，不是平均时延、P95 或性能承诺。

\begin{table}[htbp]\centering
\caption{五项真实查询的已存结果}
\begin{tabularx}{\textwidth}{@{}L{37mm}Y L{19mm}@{}}\toprule
任务 & 结果要点 & 单次耗时\\\midrule
2026 秋数学分析开课 & 返回 A1/B1 开课记录、教师和学期 & 10.92 秒\\
2026 级计算机普通主修 & 返回对应方案及总学分 164；绑定方案口径 & 12.06 秒\\
最近三条校园通知 & 返回三条通知、日期与来源链接 & 12.02 秒\\
工作日东区至西区上午校车 & 返回相应方向与时段班次 & 17.04 秒\\
二课英语角活动 & 返回 English Corner 活动、时间和报名状态 & 29.07 秒\\\bottomrule
\end{tabularx}
\end{table}

上述 5 项记录均通过对应检查。由于样本来自演示任务，无法代表所有校园问题，也不能用这些耗时比较 Flash 与 Pro 的优劣。远端内容会更新，复现时应检查字段、来源与能力选择，而非要求模型逐字输出相同文本。

\subsection{自适应启用与入站排练}
固定案例由四个模拟发言者逐步讨论数学分析选课，包含讲解、证明训练、课业负担和候选筛选等不同需求。前 19 条没有启用评课，第 20 条形成明确请求后，记录启用 \code{icourse.read} 并查询 \code{icourse.public-query.v2}。

原生入站排练记录包含 20 次状态变化、1 次自适应启用、1 份回答、1 次输出调用和 0 次 Memory 写入，最终处置为 \code{canary_completed}。台词注入接收端，不能据此声称四位成员的账号都经 QQ 网络真实发送了 20 条消息。该案例验证的是既定场景下的事件转换、参与、启用、查询与回复路径，长期自然群的误启用率尚未建立。

\subsection{广覆盖预览中暴露的问题}
2026 年 9 月 4—5 日的群聊 100 题材料包含 93 次真实模型合成历史预览和 7 项隔离故障或策略测试。真实预览终态为 89 次 response、4 次 failed；语义审查则为 60 通过、25 部分、8 失败。加上隔离测试登记，合计为 63 通过、29 部分、8 失败。不同层级登记不能写成生产端到端通过率。

问题涉及引用理解、精确格式、查询前澄清、来源缺失、摘要时间范围与发言归属，以及生成失败时正文为空等。后续录制和专项回归修复了其中若干路径，但没有与这批广覆盖预览同范围、同方法的全量重测证据。因此，新版仍保留这组结果，不以“5 项通过”覆盖较早的广泛缺陷。

另有 8 月的评课 75 题审校、二课 75 条消息模拟及 9 月 3 日固定消息报告。它们分别使用过不同模型、脚本语义或旁路审校，适合作为功能演进与回归资料，不能混算成同一 DeepSeek 评测集。

\subsection{组件与界面验证}
录制操作单记载：26/26 MCP 原子查询、5/5 自然语言问答、多个 Python 回归集合，以及 Web 114 项前端和 121 项服务端测试通过。独立 Web 试用修复报告还记录了桌面、手机与键盘场景。不同批次之间可能包含重复测试，因此本文不将这些数量累加为独立测试总数。

Memory 检索与 Bandit 评测均为合成离线数据。它们支持组件行为与计算正确性，不支持长期记忆效用、真实社交反馈或在线路由收益。当前还没有经过设计的用户研究来量化学习效率提升或管理成本下降。

\subsection{当前完成度与后续验证重点}
\begin{tabularx}{\textwidth}{@{}L{34mm}Y@{}}\toprule
阶段 & 能力与验证重点\\\midrule
当前运行能力 & QQ/预览入口、近期上下文、混合感知、确定性社交决策、按群自动参与、授权自适应启用、六类校园查询、模型档位与 Web 管理。\\
独立组件或离线成果 & Memory v2、主动调度与摘要模块、离线 Bandit、扩展 Persona 与可选服务资产。需要进一步接入和真实体验验证。\\
仍需建设 & 多步 Skill、长期群友印象、多 OC 切换、自然群参与质量、摘要时间与归属、查询延迟分解、当前广覆盖真实模型复测。\\\bottomrule
\end{tabularx}
\source{[S10, S11, S12, S13, S16]。数字直接取自已存 JSON、数据卡与结果报告；本轮离线复测单独留存。}
